%! Interpreting p-Values — Unit 6, Lesson 6.5 — Warm-Up (Answer Key)
\documentclass[10pt]{article}
\usepackage{apstats-article}
\usepackage{apstats-key}   % pulls in -boxes; do NOT also load -boxes
\newcommand{\TallMath}[1]{$\displaystyle #1\rule[-1.4em]{0pt}{3.2em}$}

\begin{document}

\pageheader{Unit 6, Lesson 6.5}{Warm-Up --- Answer Key}
\namedateperiod

\vspace{0.15in}

A quality control engineer claims that 10\% of widgets from a production line are defective.
A random sample of 150 widgets finds 22 defective.

\begin{enumerate}
  \item $H_0$: \ans{$p = 0.10$} \quad $H_a$: \ans{$p \ne 0.10$}

  \item Conditions:
        \begin{itemize}
          \item Random: \ans{SRS stated --- met}
          \item Large Counts: \ans{$150(0.10)=15\ge10$; $150(0.90)=135\ge10$ --- met}
          \item 10\%: \ans{$150 \le 0.10\times$(total production) --- met}
        \end{itemize}

  \item ($\hat p = $ \ans{$22/150\approx0.147$})

        $z = \dfrac{\ans{0.147} - \ans{0.10}}{\sqrt{\ans{0.10(0.90)}/\ans{150}}} = \dfrac{0.047}{0.02449} \approx$ \ans{$1.92$}

  \item \ansline{Both tails --- $H_a: p \ne 0.10$ is two-sided, so values as extreme as $z=1.92$ in \emph{either} direction count as evidence against $H_0$.}
\end{enumerate}

\end{document}
